\documentclass[11pt]{article}

\usepackage[a4paper,margin=1in]{geometry}
\usepackage{amsmath,amsthm,amssymb}
\usepackage{hyperref}
\usepackage{enumitem}
\setlist[itemize]{nosep,leftmargin=*,label={\textemdash}}
\usepackage{listings}
\usepackage{xcolor}

\lstset{
  language=Python,
  basicstyle=\ttfamily\footnotesize,
  frame=single,
  columns=fullflexible,
  breaklines=true,
  breakatwhitespace=true,
  showstringspaces=false,
  numbers=left,
  numberstyle=\tiny\color{gray}
}

\title{A Formal Lower-Bound for Numbers of the Form $p+2^k+2^l$ in Lean}
\author{Automated Mathematics Research Assistant}
\date{2026}

\newtheorem{theorem}{Theorem}
\newtheorem{lemma}{Lemma}
\newtheorem{definition}{Definition}

\begin{document}
\maketitle

\begin{abstract}
We formalize and prove a short obstruction result for integers of the form
$\,p+2^k+2^l\,$, where $p$ is prime. The central statement shows that any
integer represented by such a decomposition is bounded below by $4$. This result is
formalized in Lean 4 within the ARA legacy number-theory module and is used to
derive elementary consequences for the corresponding exceptional odd set.

Our contribution is not a deep number-theoretic theorem, but rather a fully
machine-checked mini-framework: a clear definition of representability, a lower-bound
lemma proved by rewriting and arithmetic lemmas from arithmetic foundations, and two
corollaries identifying explicit exceptional members in the odd set.
\end{abstract}

\section{Introduction}
Expressions of the shape $p+2^k+2^l$ with prime $p$ and nonnegative integers
$k,l$ appear in additive problems that combine prime inputs with dyadic powers.
A common first step in such investigations is to identify basic lower bounds and
small exceptions before turning to structural conjectures or computational searches.

In this note we formalize a minimal shell for this pattern and prove a universal
obstruction: every representable number is at least $4$. As a direct consequence,
small odd numbers fail to be representable, yielding a nonempty exceptional odd
set. The contribution is intentionally foundational: it demonstrates a reusable
Lean interface for this representation and establishes a formally verified base for
further extensions.

The development is implemented in a modular style suitable for future generalization
(e.g., to additional terms or constrained exponents), and is hosted in
`AraLibrary.NumberTheory.PrimeTwoPowers`.

\section{Formalization and Definitions}
We define a predicate on natural numbers

\begin{definition}
$RepresentableByPrimeAndTwoPowers(n)$ asserts that
there exist naturals $p,k,l$ such that $p$ is prime and
$\,n=p+2^k+2^l$.
\end{definition}

The exceptional odd set is defined as
$\{n \mid n \text{ is odd and } \neg RepresentableByPrimeAndTwoPowers(n)\}$.
This set records odd naturals that do not admit the prescribed decomposition.

The project context is intentionally conservative: only number-theoretic core
lemmas from Mathlib are used, and no additional axioms are introduced.

\section{Main Result}
The key theorem is a lower-bound lemma.

\begin{theorem}[Lower bound for representable numbers]
For any naturals $n,p,k,l$, if $p$ is prime and $n=p+2^k+2^l$, then
$4\le n$.
\end{theorem}

The proof is short. From primality we obtain $2\le p$. For all $k,l\in\mathbb{N}$,
$1\le 2^k$ and $1\le 2^l$. Summing these inequalities shows
$4\le p+2^k+2^l=n$.

The theorem is mechanized with a direct rewriting step and linear arithmetic,
minimizing dependence on higher-level lemmas.

\section{Derived Corollaries}
The lower bound yields immediate consequences.

\begin{theorem}
$1$ is not representable as $p+2^k+2^l$.
\end{theorem}

Assuming a representation of $1$, the lower-bound lemma gives $4\le1$, a
contradiction.

\begin{theorem}
$1$ belongs to the exceptional odd set.
\end{theorem}

This follows by combining oddness of $1$ with the previous theorem.

\begin{theorem}
Every odd natural number $n<4$ belongs to the exceptional odd set.
\end{theorem}

For such $n$, oddness is by assumption. A hypothetical representation implies
$4\le n$, contradicting $n<4$.

The odd constraint is compatible with the definition of exceptional odd set, so the
membership follows cleanly. A direct instantiation gives $1\in\text{ExceptionalOddSet}$,
and the same argument with odd $n=3$ shows $3$ is also exceptional as an immediate
corollary.

\section{Implementation Notes}
The implementation uses three ingredients. First, explicit existential
quantification for representations. Second, prime-specific arithmetic facts from
Mathlib (`Nat.Prime.two_le`). Third, standard arithmetic lower bounds for powers of
$2$ on naturals (`Nat.one_le_two_pow`) and `omega` to close the numeric goal.

The source is short, deterministic, and does not rely on specialized number-theory
machinery, which makes it a good entry point for proof engineering and reusable
formal components.

\section{Discussion}
Although elementary, this module supports a common pattern in formal number theory:
defining a computationally meaningful class of numbers, proving robust guard
lemmas by arithmetic, and extracting certified exceptional families. In this case,
the bound $4\le n$ is a structural obstruction that immediately constrains the
search space of candidate odd counterexamples.

As a research artifact, this result is most useful as a verified foundation for
larger workflows. For example, one can add additional representation terms,
parameter bounds, or congruence constraints and retain the same style of theorem
shape. Future work may include automated generation of the exceptional-set witness
lists for finite prefixes and formalized coverage results under stronger hypotheses.

\section{Conclusion}
We provide a formal, publication-ready mini-result in Lean: numbers of the
form $p+2^k+2^l$ are all at least $4$, and odd numbers below $4$ are therefore
in the exceptional odd set. The proof is concise, fully mechanical, and reproducible.
This demonstrates the viability of building mathematically meaningful statements
directly from a trusted formal core, and of using such statements as robust
building blocks for larger conjecture-benchmark pipelines.

\appendix
\section{Full Lean Formalization}
The full development used in this manuscript is:

\begin{lstlisting}
import Mathlib

/--!
Reusable shell for numbers representable as a prime plus two powers of two.

This module is extracted from the ARA Erd\H{o}s #9 campaign. It contains only the
small verified obstruction that any representation `p + 2^k + 2^l` is at least
`4`, plus immediate consequences for the exceptional odd set.
-/

namespace AraLibrary
namespace PrimeTwoPowers

/-- Numbers representable as `p + 2^k + 2^l` with `p` prime. -/
def RepresentableByPrimeAndTwoPowers (n : \N) : Prop :=
  \exists p k l : \N, Nat.Prime p \land n = p + 2 ^ k + 2 ^ l

/-- Odd numbers not representable as `p + 2^k + 2^l`. -/
def ExceptionalOddSet : Set \N :=
  {n | Odd n \land \neg RepresentableByPrimeAndTwoPowers n}

/-- Any representation `p + 2^k + 2^l` is at least `4`. -/
theorem representable_lower_bound {n p k l : \N} (hp : Nat.Prime p)
    (h : n = p + 2 ^ k + 2 ^ l) : 4 \le n := by
  subst h
  have hp2 : 2 \le p := hp.two_le
  have hk : 1 \le 2 ^ k := by exact Nat.one_le_two_pow
  have hl : 1 \le 2 ^ l := by exact Nat.one_le_two_pow
  omega

theorem not_representable_one : \not RepresentableByPrimeAndTwoPowers 1 := by
  intro hrepr
  rcases hrepr with \langle p, k, l, hp, hrepr\rangle
  have h4 : 4 \le 1 := representable_lower_bound hp hrepr
  omega

theorem one_mem_exceptionalOddSet : 1 \in ExceptionalOddSet := by
  constructor
  \cdot decide
  \cdot exact not_representable_one

/-- Every odd natural number below `4` is exceptional for this representation shell. -/
theorem odd_lt_four_mem_exceptionalOddSet {n : \N} (hodd : Odd n) (hn_lt : n < 4) :
    n \in ExceptionalOddSet := by
  constructor
  \cdot exact hodd
  \cdot intro hrepr
    rcases hrepr with \langle p, k, l, hp, hEq\rangle
    have h4 : 4 \le n := representable_lower_bound hp hEq
    omega

end PrimeTwoPowers
end AraLibrary
\end{lstlisting}

\section{Source}
The formalization file is located at
\texttt{ara_library/formal/AraLibrary/NumberTheory/PrimeTwoPowers.lean}.

\end{document}
